\subsection{齐次集合}

定义 6. 设 G 为一个群。G 在集合 E 上的作用称为传递的，如果存在元素 \(x \in E\) ，其轨道为 E。若 G 在 E 上的作用是传递的，则称 G-集 E 为齐次的。

也称 G 在 E 上传递地作用；或者说 E 是 G 下的齐次集。这等价于说 E 非空，且对 E 的任意元素 x 和 y，存在元素 \(\alpha \in G\) 使得 \(\alpha.x = y\) 。

例。若 E 是一个 G-集，则 E 的每个轨道（带诱导作用）都是 G 下的齐次集。

设 G 是一个群，H 是 G 的一个子群。考虑模 H 的左陪集集合 G/H。群 G 依法则 \((g, xH) \mapsto gxH\) 在 G/H 上左作用。设 N 为 H 的正规化子。群 N 依法则 \((xH, n) \mapsto xHn = xnH\) 在 G/H 上右作用。此作用在 H 上诱导出平凡作用，从而在取商后，N/H 在 G/H 上右作用。设 \(\phi: (N/H)^{0} \to S_{G/H}\) 为与此作用对应的同态。

命题 5. 在上述记号下，G/H 是一个齐次 G-集。映射 \(\phi\) 诱导出 \((\mathrm{N} / \mathrm{H})^0\) 到 G-集 G/H 的自同构群上的同构。

元素 \(\dot{e}=H\) 在 G/H 中的轨道是 G/H，由此得第一个断言。我们现在证明第二个断言。若 \(n\in N\) 通过右平移在 G/H 上定义的是恒等映射，则 \(\dot{e}.n=\dot{e}\) ，即 H.n=H，从而 \(n\in H\) 。因此

N/H 在 G/H 上右侧忠实作用，且 \(\phi\) 是单射。G 在 G/H 上的左作用与 N/H 在 G/H 上的右作用可交换，因此 N/H 的算子给出 G/H 到自身的 G-态射；由于它们是双射，这些态射必然是 G-自同构。于是 \(\phi\) 取值于 G/H 的 G-自同构所成的群 \(\Phi\) 。我们证明 \(\phi\) 的像就是 \(\Phi\) 。设 \(f \in \Phi\) 。通过结构转移，\(f(\dot{e})\) 在 G 中的稳定化子等于 \(\dot{e}\) 的稳定化子，即等于 H。设 \(n \in G\) 使得 \(f(\dot{e}) = n\dot{e}\) 。\(n\dot{e}\) 在 G 中的稳定化子是 \(nHn^{-1}\) （第 2 号，命题 2），因此 \(nHn^{-1} = H\) ，即 \(n \in N\) 。对 G/H 的每个元素 xH，有 \(f(xH) = f(x.\dot{e}) = x.f(\dot{e}) = xnH = xHn\) ，故 f 与由 n 定义的映射一致。

注记.(1) 设 G 是一个群，H 是 G 的一个子群，且 \(\phi: \mathbf{G} \to \mathfrak{S}_{\mathbf{G}/\mathbf{H}}\) 是对应于 G 在 G/H 上作用的同态。\(\phi\) 的核是 H 的诸共轭子群的交集（第 2 号，命题 2）。它也是包含在 H 中的最大正规子群（第 3 号）。特别地，G 忠实作用于 G/H 当且仅当 H 的诸共轭子群的交集缩减为 \(\{e\}\) 。

\begin{enumerate}
\def\labelenumi{(\arabic{enumi})}
\setcounter{enumi}{1}
\tightlist
\item
  设 G 为一个群，H 与 K 为其子群，且 H 是 K 的正规子群。则 K/H 在 G-集 G/H 上右作用，且 G/H 到 G/K 的典范映射在取商后定义出 G-集同构 \((G/H)/(K/H) \rightarrow G/K\) （参见第 4 号）。
\end{enumerate}

命题 6. 设 G 是一个群，E 是一个齐次 G-集， \(a \in E\) ，H 是 \(a\) 的稳定化子，且 K 是 G 中包含于 H 的子群。存在唯一的 G-态射 \(f\) ：G/K 到 E，使得 \(f(e.K) = a\) 。若 K = H，则 \(f\) 是同构。

若 \(f\) 是一个解，则 \(f(x, \mathbf{K}) = x.a\) 对所有 \(x\) 属于 \(\mathbf{G}\) 者成立，由此得唯一性；我们证明存在性。由 \(a\) 定义的轨道映射与等价关系 \(y \in x\mathbf{K}\) 在 \(\mathbf{G}\) 上相容。因为，若 \(y = xk, k \in \mathbf{K}\) ，则

\[
y. a = x k. a = x. a.
\]

映射 \(f\) 便这样由 \(G / K\) 导出到 \(H\) ，它满足 \(f(x, \mathbf{K}) = x.a\) 对所有 \(x\) 属于 \(G\) 成立。此映射是一个 G-态射，且 \(f(\mathbf{K}) = a\) 。此映射是满射的，因为它的像是 \(E\) 的一个非空稳定子集。现在设 \(K = H\) ，我们证明 \(f\) 是单射。若 \(f(x, \mathbf{H}) = f(y, \mathbf{H})\) ，则 \(x.a = y.a\) ，从而 \(x^{-1}y.a = a\) ，即 \(x^{-1}y \in H\) ，于是 \(x.H = y.H\) 。

定理 1. 设 G 是一个群。

\begin{enumerate}
\def\labelenumi{(\alph{enumi})}
\item
  每个齐次 G-集都同构于形如 G/H 的齐次 G-集，其中 H 是 G 的一个子群。
\item
  设 H 和 H' 是 G 的两个子群。G-集 G/H 与 G/H' 同构当且仅当 H 与 H' 共轭。
\end{enumerate}

由于齐次 G-集非空，断言 (a) 由命题 6 得出。我们证明 (b)。设 \(f: \mathrm{G} / \mathrm{H} \to \mathrm{G} / \mathrm{H}'\) 是一个 G-集同构。子群 H 是 H 的稳定化子，因此通过结构转移，它是 G/H' 中某元素的稳定化子。于是子群 H 与 H' 共轭（第 2 号，命题 2）。若 H' = αHα⁻¹，则 H' 是元素 α. H 在 G/H 中的稳定化子（第 2 号，命题 2），因此 G/H' 同构于 G/H（命题 6）。

例子。(1) 设 E 是一个非空集合。群 \(S_{E}\) 在 E 上传递地作用。若 x 和 y 是 E 的两个元素，则映射 \(\tau: E \to E\) （ \(\tau(x) = y\) , \(\tau(y) = x\) ，且 \(\tau(z) = z\) 对 \(z \neq x\) , y 成立）是 E 的一个置换。设 \(a \in E\) 。a 的稳定化子等同于 \(S_{F}\) ，其中 \(F = E - \{a\}\) 。于是齐次 G-集 E 同构于 \(S_{E}/S_{F}\) 。

\begin{enumerate}
\def\labelenumi{(\arabic{enumi})}
\setcounter{enumi}{1}
\tightlist
\item
  设 \(\mathbf{E}\) 是含 \(n\) 个元素的集合，\((p_i)_{i\in \mathbf{I}}\) 是整数 \(>0\) 的有限族，使得 \(\sum_{i}p_{i} = n\) 。设 \(\mathbf{X}\) 是形如 \((\mathrm{F}_i)_{i\in \mathbf{I}}\) 的 \(\mathbf{E}\) 的分拆中满足 \(\operatorname {Card}(\mathrm{F}_i) = p_i\) （对所有 \(i\) ）者所成的集合。群 \(\mathfrak{S}_{\mathbf{E}}\) 在 \(\mathbf{X}\) 上传递地作用。元素 \((\mathrm{F}_i)_{i\in \mathbf{I}}\) （属于 \(\mathbf{X}\) ）的稳定化子 H 典范同构于 \(\prod_{i\in \mathbf{I}}\mathfrak{S}_{\mathbf{F}_i}\) ，从而其阶为 \(\prod_{i\in \mathbf{I}}p_i!\) 。应用定理 1 与 §4，第 4 号，命题 4 的推论，可得到下述事实的一个新证明：
\end{enumerate}

\[
\operatorname{Card} (\mathbf {X}) = \frac {n !}{\prod_ {i \in I} p _ {i} !}.
\]

特别地，取 \(\mathbf{I} = \{1,2,\dots,r\}\) , \(\mathbf{E} = \{1,2,\dots,n\}\) ,

\[
\mathrm{F} _ {i} = \left\{p _ {1} + \dots + p _ {i - 1} + 1, \dots , p _ {1} + \dots + p _ {i} \right\}
\]

对 \(1 \leqslant i \leqslant r\) ，设 \(S\) 为 \(\tau \in \mathfrak{S}_{E}\) 中使 \(\tau|_{F_i}\) 对 \(1 \leqslant i \leqslant r\) 递增者所成的集合。若 \((G_1, \ldots, G_r) \in X\) ，则存在唯一的 \(\tau \in S\) 把 \((F_1, \ldots, F_r)\) 映为 \((G_1, \ldots, G_r)\) 。换言之，\(\mathfrak{S}_{E}\) 模 \(H\) 的每个左陪集与 \(S\) 恰好交于一点。

\begin{enumerate}
\def\labelenumi{(\arabic{enumi})}
\setcounter{enumi}{2}
\tightlist
\item
  *设 \(n\) 为整数 \(\geqslant 1\) 。正交群 \(\mathbf{O}(n, \mathbf{R})\) 在 \(\mathbf{S}_{n-1}\) （ \(\mathbf{R}^n\) 中的单位球面）上传递地作用。点 \((0, \ldots, 0, 1)\) 的稳定化子等同于正交群 \(\mathbf{O}(n-1, \mathbf{R})\) 。于是齐次 \(\mathbf{O}(n, \mathbf{R})\) -集 \(\mathbf{S}_{n-1}\) 同构于 \(\mathbf{O}(n, \mathbf{R}) / \mathbf{O}(n-1, \mathbf{R})\) 。*
\end{enumerate}

